% ============================================================================
% COURS 04 — EXERCICES DE CINÉMATIQUE — PARTIE A
% Source : 5-Cinematique2022.docx
% ============================================================================

\section{Exercices du chapitre}
\label{sec:cine-exercices}

\subsection{Lecture d'une loi en trapèze}

Le profil d'accélération d'un mobile est défini par morceaux entre
$0$ et $3{,}2\ \mathrm s$. La position et la vitesse initiales sont
nulles.

\TLCineImage[.76\linewidth]{images/image102.png}

\begin{enumerate}
  \item Tracer la vitesse en intégrant graphiquement l'accélération.
  \item Tracer la position en calculant les aires sous la courbe de vitesse.
  \item Indiquer les valeurs maximales de vitesse et de déplacement.
  \item Vérifier la continuité de $x(t)$ et $v(t)$ aux changements de phase.
\end{enumerate}

\subsection{Maison DÔME}

Une maison circulaire de diamètre $16\ \mathrm m$ s'oriente afin de
suivre ou d'éviter le soleil. Lors de la coupure du moteur, sa vitesse
de rotation décroît linéairement jusqu'à zéro en $T_a=8\ \mathrm s$.
La fréquence initiale vaut $30\times10^{-3}\ \mathrm{tr\,min^{-1}}$.

\TLCineImage[.48\linewidth]{images/image103.pdf}

\begin{enumerate}
  \item Expliquer pourquoi la maison continue de tourner après la coupure.
  \item Calculer l'accélération angulaire en $\mathrm{rad\,s^{-2}}$.
  \item Déterminer l'angle parcouru pendant l'arrêt.
  \item Calculer la distance parcourue au niveau de la porte, située à $8\ \mathrm m$ de l'axe.
  \item Vérifier le critère de distance d'arrêt inférieure à $15\ \mathrm{cm}$.
\end{enumerate}

\subsection{Véhicule électrique F-City}

La F-City atteint $30\ \mathrm{km\,h^{-1}}$ en $5{,}5\ \mathrm s$, sa
vitesse maximale vaut $65\ \mathrm{km\,h^{-1}}$ et son autonomie
minimale est $80\ \mathrm{km}$.

\TLCineImage[.52\linewidth]{images/image104.pdf}

\begin{enumerate}
  \item Calculer l'accélération moyenne entre 0 et $30\ \mathrm{km\,h^{-1}}$.
  \item Calculer la durée maximale d'un trajet de $80\ \mathrm{km}$ à vitesse maximale.
  \item À partir du profil de vitesse fourni, déterminer les distances de chaque phase.
  \item Calculer la distance totale du parcours et vérifier l'autonomie.
\end{enumerate}

\subsection{Éolienne à pales orientables}

Le rotor 1 tourne autour de $(O,\vec x_0)$ avec l'angle $\theta$. Une
pale 2 pivote par rapport au rotor avec l'angle de calage $\alpha$.
On donne :
\[
\overrightarrow{OA}=a\vec y_1,
\qquad
\overrightarrow{AG}=x\vec x_2+y\vec y_2+z\vec z_2.
\]

\TLCineImage[.48\linewidth]{images/image105.jpeg}

\begin{enumerate}
  \item Tracer les deux figures de changement de base et donner les vecteurs rotation.
  \item Déterminer $\vec V(A\in1/0)$.
  \item Déterminer $\vec V(G\in2/0)$ par composition.
  \item Déterminer $\vec A(A\in1/0)$ puis $\vec A(G\in2/0)$ par dérivation.
  \item Identifier les termes centrifuge, tangentiel et de Coriolis.
\end{enumerate}

\subsection{Centrifugeuse d'entraînement}

Le bras 1 tourne autour de l'axe vertical $(O,\vec z)$ avec l'angle
$\alpha$. La cabine 2 pivote autour de $(A,\vec y_1)$ avec l'angle
$\beta$. On donne $\overrightarrow{OA}=a\vec x_1$ et
$\overrightarrow{AG}=b\vec z_2$.

\TLCineImage[.58\linewidth]{images/image106.png}

\begin{enumerate}
  \item Tracer les figures de changement de base et les vecteurs rotation.
  \item Déterminer $\vec\Omega_{2/0}$.
  \item Déterminer $\vec V(G\in2/0)$ par composition.
  \item Déterminer $\vec A(G\in2/0)$ par dérivation du résultat précédent.
  \item Interpréter les composantes de l'accélération ressentie par le pilote.
\end{enumerate}

\subsection{Robot manipulateur à deux axes}

Le corps 2 tourne par rapport au socle 1 autour de
$(O_1,\vec z_1)$ avec l'angle $\varphi$. Le bras 3 coulisse suivant
$\vec y_2=\vec y_3$ avec
$\overrightarrow{AG_3}=y(t)\vec y_3$ et
$\overrightarrow{O_1A}=a\vec x_2$.

\TLCineImage[.50\linewidth]{images/image109.jpeg}

\begin{enumerate}
  \item Donner la figure de changement de base et $\vec\Omega_{2/1}$.
  \item Écrire le torseur de la glissière 3/2.
  \item Déterminer $\vec V(G_3\in3/1)$.
  \item Déterminer $\vec A(G_3\in3/1)$.
  \item Repérer le terme de Coriolis lorsque $\dot y\neq0$ et $\dot\varphi\neq0$.
\end{enumerate}

\subsection{Transporteur horizontal}

Un chariot 2 se déplace suivant $\vec x_1$ avec la coordonnée $x(t)$.
Un bras 3 de longueur $L_3$ est articulé en $B$ sur le chariot ; son
centre de gravité est tel que
$\overrightarrow{BG_3}=L_3\vec y_3$ et son orientation est $\theta$.

\TLCineImage[.50\linewidth]{images/image110.jpeg}

\begin{enumerate}
  \item Donner les torseurs $\{\mathcal V_{2/1}\}_A$ et $\{\mathcal V_{3/2}\}_A$.
  \item Déterminer $\vec A(A\in2/1)$.
  \item Déterminer $\vec V(G_3\in3/1)$ puis $\vec A(G_3\in3/1)$.
  \item Exprimer les résultats en fonction de $\dot x$, $\ddot x$, $\dot\theta$, $\ddot\theta$ et $L_3$.
\end{enumerate}

\subsection{Centre d'usinage cinq axes HSM 600U}

Le centre possède trois translations $x(t)$, $y(t)$, $z(t)$ et deux
rotations $\alpha(t)$ et $\beta(t)$. Les axes linéaires atteignent
$40\ \mathrm{m\,min^{-1}}$ ; les axes rotatifs atteignent respectivement
$150$ et $250\ \mathrm{tr\,min^{-1}}$.

\begin{center}
\TLCineImage[.29\linewidth]{images/image111.jpeg}\hfill
\TLCineImage[.31\linewidth]{images/image112.png}\hfill
\TLCineImage[.34\linewidth]{images/image113.png}
\end{center}

\begin{enumerate}
  \item Identifier la nature des cinq mouvements d'axes.
  \item Déterminer les relations entre les bases des solides en translation.
  \item Écrire le vecteur position de l'extrémité d'outil $P$ dans $R_0$.
  \item Déterminer $\vec V(P\in5/0)$ et sa norme maximale.
  \item Déterminer la position, la vitesse et l'accélération du point $O_3$ porté par la table orientable.
  \item Retrouver la vitesse par composition des torseurs cinématiques.
  \item Écrire la vitesse relative de l'outil par rapport à un point de la pièce usinée.
\end{enumerate}

\subsection{Échelle pivotante automatique EPAS}

Le berceau 2 tourne par rapport à la tourelle 1 avec l'angle $\theta$.
L'échelle 3 se déploie suivant $\vec x_2$ avec la course $\lambda(t)$.
Pendant le dressage, la tourelle reste fixe.

\TLCineImage[.65\linewidth]{images/image172.png}

\begin{enumerate}
  \item Donner la nature des mouvements 2/0, 3/2 et 3/0.
  \item Décrire les trajectoires d'un point de l'échelle dans ces trois mouvements.
  \item Déterminer $\vec\Omega_{3/0}$.
  \item Déterminer la vitesse d'un point $G$ de la nacelle.
  \item Déterminer son accélération et identifier les termes liés à $\lambda$, $\theta$ et leurs dérivées.
  \item Vérifier le critère d'accélération maximale du cahier des charges.
\end{enumerate}
